% source: https://tex.stackexchange.com/a/765262 (question 726438, score 8, block 2)
\documentclass[border=5pt,tikz]{standalone}% compile with lualatex only
\usepackage[svgnames]{xcolor}
\usepackage[3d]{luadraw}%https://github.com/pfradin/luadraw
\usepackage{fourier-otf}
%https://tex.stackexchange.com/questions/726438/drawing-the-steinmetz-solid-with-tikz
\begin{document}
\begin{luadraw}{name=Steinmetz_solid3}
local ld = luadraw
local pt3d = ld.pt3d
local Origin, vecI, vecJ, vecK, M = pt3d.Origin, pt3d.vecI, pt3d.vecJ, pt3d.vecK, pt3d.M
local sin, cos, pi = math.sin, math.cos, math.pi
local g = ld.graph3d:new{window={-3,3,-3,3},size={12,12}, 
    viewdir={25,50}, bg="gray", bbox=false}
local r = 2
local axe = {Origin, vecK}
local p = function(u,v) 
    local M1, M2 = M(cos(u),cos(u),sin(u)), M(cos(u),-cos(u),sin(u))
    return r*((1-v)*M1+v*M2)
end
local S1 = ld.surface(p,-pi/2,pi/2,0,1,{20,2}) --first quarter for the blue cylinder 
table.append(S1, ld.rotate3d(S1,180,axe) ) -- addition of the opposite quarter
local L = ld.border(S1) -- the points common to the two cylinders
local S2 = ld.rotate3d(S1, 90, axe) -- first and second quarters of the red cylinder

g:Dscene3d( 
    g:addFacet(S1, {color="white",contrast=0, edge=true, edgecolor="blue"}),
    g:addFacet(S2, {color="white",contrast=0, edge=true, edgecolor="red"}),
    g:addPolyline(L, {color="Magenta", width=18, hidden=true})
    )
g:Show()
\end{luadraw}
\end{document}
